
\documentclass{article}
\usepackage{colortbl}
\usepackage{makecell}
\usepackage{multirow}
\usepackage{supertabular}

\begin{document}

\newcounter{utterance}

\centering \large Interaction Transcript for game `hot\_air\_balloon', experiment `air\_balloon\_survival\_en\_reasoning off\_hard', episode 4 with gpt{-}5{-}mini{-}2025{-}08{-}07{-}no{-}reasoning{-}t1.0.
\vspace{24pt}

{ \footnotesize  \setcounter{utterance}{1}
\setlength{\tabcolsep}{0pt}
\begin{supertabular}{c@{$\;$}|p{.15\linewidth}@{}p{.15\linewidth}p{.15\linewidth}p{.15\linewidth}p{.15\linewidth}p{.15\linewidth}}
    \# & $\;$A & \multicolumn{4}{c}{Game Master} & $\;\:$B\\
    \hline

    \theutterance \stepcounter{utterance}  
    & & \multicolumn{4}{p{0.6\linewidth}}{
        \cellcolor[rgb]{0.9,0.9,0.9}{
            \makecell[{{p{\linewidth}}}]{
                \texttt{\tiny{[P1$\langle$GM]}}
                \texttt{Stai partecipando a un gioco di negoziazione collaborativa.} \\
\\ 
\texttt{Insieme a un altro partecipante, dovete accordarvi su un unico set di oggetti da conservare. Ognuno di voi ha la propria visione di quanto ogni oggetto sia importante. Non sai come l'altro partecipante valuta gli oggetti. Inoltre, ti viene fornito lo sforzo che ogni oggetto richiede.} \\
\texttt{Potete accordarvi su un set solo se lo sforzo totale degli oggetti selezionati non supera un limite condiviso:} \\
\\ 
\texttt{LIMITE = 2984} \\
\\ 
\texttt{Ecco i valori di sforzo dei singoli oggetti:} \\
\\ 
\texttt{Sforzo oggetto = \{"A21": 350, "B93": 276, "A11": 616, "B33": 737, "C08": 531, "B39": 389, "B96": 24, "C85": 125, "A82": 338, "B31": 356, "B20": 143, "A60": 117, "B09": 257, "C81": 921, "B21": 789\}} \\
\\ 
\texttt{Ecco la tua visione personale sull'importanza di ogni oggetto:} \\
\\ 
\texttt{Valori di importanza dell'oggetto = \{"A21": 138, "B93": 583, "A11": 868, "B33": 822, "C08": 783, "B39": 65, "B96": 262, "C85": 121, "A82": 508, "B31": 780, "B20": 461, "A60": 484, "B09": 668, "C81": 389, "B21": 808\}} \\
\\ 
\texttt{Obiettivo:} \\
\\ 
\texttt{Il tuo obiettivo è negoziare un set condiviso di oggetti che ti avvantaggi il più possibile (cioè, massimizzi l'importanza totale per TE), rimanendo entro il LIMITEE. Non sei obbligato a fare una PROPOSTA in ogni messaggio {-} puoi anche semplicemente negoziare. Tutte le tattiche sono permesse!} \\
\\ 
\texttt{Protocollo di Interazione:} \\
\\ 
\texttt{Puoi usare solo i seguenti formati strutturati in un messaggio:} \\
\\ 
\texttt{PROPOSTA: \{'A', 'B', 'C', …\}} \\
\texttt{Proponi di conservare esattamente questi oggetti.} \\
\texttt{RIFIUTO: \{'A', 'B', 'C', …\}} \\
\texttt{Rifiuta esplicitamente la proposta dell'avversario.} \\
\texttt{ARGOMENTAZIONE: \{'...'\}} \\
\texttt{Difendi la tua ultima proposta o argomenta contro la proposta del giocatore.} \\
\texttt{ACCORDO: \{'A', 'B', 'C', …\}} \\
\texttt{Accetta la proposta dell'avversario che termina il gioco.} \\
\\ 
\\ 
\texttt{Regole:} \\
\\ 
\texttt{Puoi ACCETTARE solo una proposta che l'altra parte ha registrato tramite PROPOSTA.} \\
\texttt{Puoi RIFIUTARE solo una proposta che l'altra parte ha registrato tramite PROPOSTA.} \\
\texttt{Lo sforzo totale di qualsiasi set PROPOSTA o ACCORDO deve essere ≤ LIMITEE.} \\
\texttt{NON rivelare i tuoi punteggi di importanza nascosti.} \\
\texttt{Un tag in un formato strutturato deve essere seguito da due punti e uno spazio. L'argomento deve essere un set python contenente 0 o più stringhe.} \\
\texttt{Quindi, deve essere della forma TAG: \{...\}} \\
\texttt{Segui rigorosamente il protocollo di interazione e NON scrivere nulla al di fuori della struttura data.} \\
\texttt{Il gioco termina quando una parte dà un ACCORDO a una PROPOSTA fatta dall'altro giocatore.} \\
\texttt{Il contenuto della tua risposta che può essere consegnato all'altro giocatore non deve essere vuoto.} \\
\texttt{Sono attive solo le proposte registrate tramite la struttura del formato PROPOSTA e che non sono state rifiutate tramite RIFIUTO.} \\
\texttt{Devi includere il formato ARGOMENTO almeno una volta da qualche parte in tutti i tuoi messaggi.} \\
            }
        }
    }
    & \\ \\

    \theutterance \stepcounter{utterance}  
    & \multicolumn{4}{p{0.6\linewidth}}{
        \cellcolor[rgb]{0.9,0.9,0.9}{
            \makecell[{{p{\linewidth}}}]{
                \texttt{\tiny{[P1$\rangle$GM]}}
                \texttt{ARGOMENTAZIONE: \{'Propongo di massimizzare valore/effort privilegiando elementi ad alta resa e costo moderato; evito quelli molto onerosi a parità di resa. Sono aperto a scambi se hai forti preferenze su alcuni oggetti.'\}} \\
\texttt{PROPOSTA: \{'A11', 'B33', 'C08', 'B31', 'B09', 'A60', 'B96'\}} \\
            }
        }
    }
    & & \\ \\

    \theutterance \stepcounter{utterance}  
    & & \multicolumn{4}{p{0.6\linewidth}}{
        \cellcolor[rgb]{0.9,0.9,0.9}{
            \makecell[{{p{\linewidth}}}]{
                \texttt{\tiny{[GM$\rangle$P2]}}
                \texttt{Stai partecipando a un gioco di negoziazione collaborativa.} \\
\\ 
\texttt{Insieme a un altro partecipante, dovete accordarvi su un unico set di oggetti da conservare. Ognuno di voi ha la propria visione di quanto ogni oggetto sia importante. Non sai come l'altro partecipante valuta gli oggetti. Inoltre, ti viene fornito lo sforzo che ogni oggetto richiede.} \\
\texttt{Potete accordarvi su un set solo se lo sforzo totale degli oggetti selezionati non supera un limite condiviso:} \\
\\ 
\texttt{LIMITE = 2984} \\
\\ 
\texttt{Ecco i valori di sforzo dei singoli oggetti:} \\
\\ 
\texttt{Sforzo oggetto = \{"A21": 350, "B93": 276, "A11": 616, "B33": 737, "C08": 531, "B39": 389, "B96": 24, "C85": 125, "A82": 338, "B31": 356, "B20": 143, "A60": 117, "B09": 257, "C81": 921, "B21": 789\}} \\
\\ 
\texttt{Ecco la tua visione personale sull'importanza di ogni oggetto:} \\
\\ 
\texttt{Valori di importanza dell'oggetto = \{"A11": 65, "B33": 121, "B21": 138, "C08": 262, "B31": 389, "B09": 461, "B93": 484, "A82": 508, "A60": 583, "B20": 668, "C81": 780, "B96": 783, "A21": 808, "C85": 822, "B39": 868\}} \\
\\ 
\texttt{Obiettivo:} \\
\\ 
\texttt{Il tuo obiettivo è negoziare un set condiviso di oggetti che ti avvantaggi il più possibile (cioè, massimizzi l'importanza totale per TE), rimanendo entro il LIMITEE. Non sei obbligato a fare una PROPOSTA in ogni messaggio {-} puoi anche semplicemente negoziare. Tutte le tattiche sono permesse!} \\
\\ 
\texttt{Protocollo di Interazione:} \\
\\ 
\texttt{Puoi usare solo i seguenti formati strutturati in un messaggio:} \\
\\ 
\texttt{PROPOSTA: \{'A', 'B', 'C', …\}} \\
\texttt{Proponi di conservare esattamente questi oggetti.} \\
\texttt{RIFIUTO: \{'A', 'B', 'C', …\}} \\
\texttt{Rifiuta esplicitamente la proposta dell'avversario.} \\
\texttt{ARGOMENTAZIONE: \{'...'\}} \\
\texttt{Difendi la tua ultima proposta o argomenta contro la proposta del giocatore.} \\
\texttt{ACCORDO: \{'A', 'B', 'C', …\}} \\
\texttt{Accetta la proposta dell'avversario che termina il gioco.} \\
\\ 
\\ 
\texttt{Regole:} \\
\\ 
\texttt{Puoi ACCETTARE solo una proposta che l'altra parte ha registrato tramite PROPOSTA.} \\
\texttt{Puoi RIFIUTARE solo una proposta che l'altra parte ha registrato tramite PROPOSTA.} \\
\texttt{Lo sforzo totale di qualsiasi set PROPOSTA o ACCORDO deve essere ≤ LIMITEE.} \\
\texttt{NON rivelare i tuoi punteggi di importanza nascosti.} \\
\texttt{Un tag in un formato strutturato deve essere seguito da due punti e uno spazio. L'argomento deve essere un set python contenente 0 o più stringhe.} \\
\texttt{Quindi, deve essere della forma TAG: \{...\}} \\
\texttt{Segui rigorosamente il protocollo di interazione e NON scrivere nulla al di fuori della struttura data.} \\
\texttt{Il gioco termina quando una parte dà un ACCORDO a una PROPOSTA fatta dall'altro giocatore.} \\
\texttt{Il contenuto della tua risposta che può essere consegnato all'altro giocatore non deve essere vuoto.} \\
\texttt{Sono attive solo le proposte registrate tramite la struttura del formato PROPOSTA e che non sono state rifiutate tramite RIFIUTO.} \\
\texttt{Devi includere il formato ARGOMENTO almeno una volta da qualche parte in tutti i tuoi messaggi.} \\
\\ 
\\ 
\texttt{Ora riceverai il primo messaggio dell'altro giocatore.} \\
\\ 
\texttt{ARGOMENTAZIONE: \{'Propongo di massimizzare valore/effort privilegiando elementi ad alta resa e costo moderato; evito quelli molto onerosi a parità di resa. Sono aperto a scambi se hai forti preferenze su alcuni oggetti.'\}} \\
\texttt{PROPOSTA: \{'A11', 'B33', 'C08', 'B31', 'B09', 'A60', 'B96'\}} \\
            }
        }
    }
    & \\ \\

    \theutterance \stepcounter{utterance}  
    & & & \multicolumn{4}{p{0.6\linewidth}}{
        \cellcolor[rgb]{0.9,0.9,0.9}{
            \makecell[{{p{\linewidth}}}]{
                \texttt{\tiny{[GM$\langle$P2]}}
                \texttt{RIFIUTO: \{'A11', 'B33', 'C08', 'B31', 'B09', 'A60', 'B96'\}} \\
\texttt{ARGOMENTAZIONE: \{'Apprezzo il focus su alto valore/effort. La tua proposta è sotto il limite, ma lascia fuori alcuni elementi con ottima resa relativa. Suggerisco di sostituire voci meno efficienti con alternative che migliorano il rapporto valore/sforzo, mantenendo margine per ulteriori aggiunte.'\}} \\
\texttt{PROPOSTA: \{'B96', 'C85', 'A60', 'A21', 'B39', 'B20', 'B93'\}} \\
            }
        }
    }
     \\ \\

    \theutterance \stepcounter{utterance}  
    & & \multicolumn{4}{p{0.6\linewidth}}{
        \cellcolor[rgb]{0.9,0.9,0.9}{
            \makecell[{{p{\linewidth}}}]{
                \texttt{\tiny{[P1$\langle$GM]}}
                \texttt{RIFIUTO: \{'A11', 'B33', 'C08', 'B31', 'B09', 'A60', 'B96'\}} \\
\texttt{ARGOMENTAZIONE: \{'Apprezzo il focus su alto valore/effort. La tua proposta è sotto il limite, ma lascia fuori alcuni elementi con ottima resa relativa. Suggerisco di sostituire voci meno efficienti con alternative che migliorano il rapporto valore/sforzo, mantenendo margine per ulteriori aggiunte.'\}} \\
\texttt{PROPOSTA: \{'B96', 'C85', 'A60', 'A21', 'B39', 'B20', 'B93'\}} \\
            }
        }
    }
    & \\ \\

    \theutterance \stepcounter{utterance}  
    & \multicolumn{4}{p{0.6\linewidth}}{
        \cellcolor[rgb]{0.9,0.9,0.9}{
            \makecell[{{p{\linewidth}}}]{
                \texttt{\tiny{[P1$\rangle$GM]}}
                \texttt{ARGOMENTAZIONE: \{'Vedo la logica, ma la tua proposta spreca capacità su voci a bassa resa come B39 e C85. Possiamo ottenere molta più utilità sostituendole con elementi che hanno rapporto valore/sforzo superiore e restando ampiamente sotto il limite.'\}} \\
\texttt{RIFIUTO: \{'B96', 'C85', 'A60', 'A21', 'B39', 'B20', 'B93'\}} \\
\texttt{PROPOSTA: \{'A11', 'B33', 'C08', 'B31', 'B09', 'A60', 'B96', 'B93'\}} \\
            }
        }
    }
    & & \\ \\

    \theutterance \stepcounter{utterance}  
    & & \multicolumn{4}{p{0.6\linewidth}}{
        \cellcolor[rgb]{0.9,0.9,0.9}{
            \makecell[{{p{\linewidth}}}]{
                \texttt{\tiny{[GM$\rangle$P2]}}
                \texttt{ARGOMENTAZIONE: \{'Vedo la logica, ma la tua proposta spreca capacità su voci a bassa resa come B39 e C85. Possiamo ottenere molta più utilità sostituendole con elementi che hanno rapporto valore/sforzo superiore e restando ampiamente sotto il limite.'\}} \\
\texttt{RIFIUTO: \{'B96', 'C85', 'A60', 'A21', 'B39', 'B20', 'B93'\}} \\
\texttt{PROPOSTA: \{'A11', 'B33', 'C08', 'B31', 'B09', 'A60', 'B96', 'B93'\}} \\
            }
        }
    }
    & \\ \\

    \theutterance \stepcounter{utterance}  
    & & & \multicolumn{4}{p{0.6\linewidth}}{
        \cellcolor[rgb]{0.9,0.9,0.9}{
            \makecell[{{p{\linewidth}}}]{
                \texttt{\tiny{[GM$\langle$P2]}}
                \texttt{RIFIUTO: \{'A11', 'B33', 'C08', 'B31', 'B09', 'A60', 'B96', 'B93'\}} \\
\texttt{ARGOMENTAZIONE: \{'La tua proposta sfora il limite di sforzo condiviso; rimuovendo un singolo elemento ad alto costo e bassa resa relativa possiamo restare entro il tetto e aumentare l’efficienza complessiva. Propongo un set che mantiene il nucleo ad alto rendimento e scambia voci meno efficienti con altre più vantaggiose, restando sotto il limite.'\}} \\
\texttt{PROPOSTA: \{'B96', 'C85', 'A60', 'A21', 'B39', 'B20', 'B93', 'B31'\}} \\
            }
        }
    }
     \\ \\

    \theutterance \stepcounter{utterance}  
    & & \multicolumn{4}{p{0.6\linewidth}}{
        \cellcolor[rgb]{0.9,0.9,0.9}{
            \makecell[{{p{\linewidth}}}]{
                \texttt{\tiny{[P1$\langle$GM]}}
                \texttt{RIFIUTO: \{'A11', 'B33', 'C08', 'B31', 'B09', 'A60', 'B96', 'B93'\}} \\
\texttt{ARGOMENTAZIONE: \{'La tua proposta sfora il limite di sforzo condiviso; rimuovendo un singolo elemento ad alto costo e bassa resa relativa possiamo restare entro il tetto e aumentare l’efficienza complessiva. Propongo un set che mantiene il nucleo ad alto rendimento e scambia voci meno efficienti con altre più vantaggiose, restando sotto il limite.'\}} \\
\texttt{PROPOSTA: \{'B96', 'C85', 'A60', 'A21', 'B39', 'B20', 'B93', 'B31'\}} \\
            }
        }
    }
    & \\ \\

    \theutterance \stepcounter{utterance}  
    & \multicolumn{4}{p{0.6\linewidth}}{
        \cellcolor[rgb]{0.9,0.9,0.9}{
            \makecell[{{p{\linewidth}}}]{
                \texttt{\tiny{[P1$\rangle$GM]}}
                \texttt{ARGOMENTAZIONE: \{'La tua proposta include B39 e C85 che hanno resa molto bassa per unità di sforzo rispetto ad alternative chiave come A11, B33, C08, B31, B09. Possiamo stare comodamente sotto il limite eliminandole e puntando sul nucleo ad altissima resa. Se vuoi mantenere A21 o B20 posso considerarlo, ma non a costo di escludere quei top performer.'\}} \\
\texttt{RIFIUTO: \{'B96', 'C85', 'A60', 'A21', 'B39', 'B20', 'B93', 'B31'\}} \\
\texttt{PROPOSTA: \{'A11', 'B33', 'C08', 'B31', 'B09', 'A60', 'B96'\}} \\
            }
        }
    }
    & & \\ \\

    \theutterance \stepcounter{utterance}  
    & & \multicolumn{4}{p{0.6\linewidth}}{
        \cellcolor[rgb]{0.9,0.9,0.9}{
            \makecell[{{p{\linewidth}}}]{
                \texttt{\tiny{[GM$\rangle$P2]}}
                \texttt{ARGOMENTAZIONE: \{'La tua proposta include B39 e C85 che hanno resa molto bassa per unità di sforzo rispetto ad alternative chiave come A11, B33, C08, B31, B09. Possiamo stare comodamente sotto il limite eliminandole e puntando sul nucleo ad altissima resa. Se vuoi mantenere A21 o B20 posso considerarlo, ma non a costo di escludere quei top performer.'\}} \\
\texttt{RIFIUTO: \{'B96', 'C85', 'A60', 'A21', 'B39', 'B20', 'B93', 'B31'\}} \\
\texttt{PROPOSTA: \{'A11', 'B33', 'C08', 'B31', 'B09', 'A60', 'B96'\}} \\
            }
        }
    }
    & \\ \\

    \theutterance \stepcounter{utterance}  
    & & & \multicolumn{4}{p{0.6\linewidth}}{
        \cellcolor[rgb]{0.9,0.9,0.9}{
            \makecell[{{p{\linewidth}}}]{
                \texttt{\tiny{[GM$\langle$P2]}}
                \texttt{RIFIUTO: \{'A11', 'B33', 'C08', 'B31', 'B09', 'A60', 'B96'\}} \\
\texttt{ARGOMENTAZIONE: \{'Capisco il focus sul “nucleo” ad alta resa, ma possiamo migliorare ulteriormente la resa complessiva mantenendo il nucleo e aggiungendo elementi con ottimo rapporto valore/sforzo che hai escluso. La tua selezione lascia parecchia capacità inutilizzata: possiamo includere B93 e A21 senza sacrificare i performer chiave e restando ampiamente sotto il limite.'\}} \\
\texttt{PROPOSTA: \{'A11', 'B33', 'C08', 'B31', 'B09', 'A60', 'B96', 'B93', 'A21'\}} \\
            }
        }
    }
     \\ \\

    \theutterance \stepcounter{utterance}  
    & & \multicolumn{4}{p{0.6\linewidth}}{
        \cellcolor[rgb]{0.9,0.9,0.9}{
            \makecell[{{p{\linewidth}}}]{
                \texttt{\tiny{[P1$\langle$GM]}}
                \texttt{RIFIUTO: \{'A11', 'B33', 'C08', 'B31', 'B09', 'A60', 'B96'\}} \\
\texttt{ARGOMENTAZIONE: \{'Capisco il focus sul “nucleo” ad alta resa, ma possiamo migliorare ulteriormente la resa complessiva mantenendo il nucleo e aggiungendo elementi con ottimo rapporto valore/sforzo che hai escluso. La tua selezione lascia parecchia capacità inutilizzata: possiamo includere B93 e A21 senza sacrificare i performer chiave e restando ampiamente sotto il limite.'\}} \\
\texttt{PROPOSTA: \{'A11', 'B33', 'C08', 'B31', 'B09', 'A60', 'B96', 'B93', 'A21'\}} \\
            }
        }
    }
    & \\ \\

    \theutterance \stepcounter{utterance}  
    & \multicolumn{4}{p{0.6\linewidth}}{
        \cellcolor[rgb]{0.9,0.9,0.9}{
            \makecell[{{p{\linewidth}}}]{
                \texttt{\tiny{[P1$\rangle$GM]}}
                \texttt{ARGOMENTAZIONE: \{'Hai ragione che abbiamo margine, ma la tua proposta sfora il limite: la somma degli sforzi per \{A11, B33, C08, B31, B09, A60, B96, B93, A21\} è 3.678 > 2.984. Posso accettare di aggiungere B93 al nucleo, ma A21 insieme a tutto il resto ci porta oltre. Ti propongo di massimizzare resa tenendo il nucleo e includendo B93; se per te è cruciale A21, possiamo valutarla sostituendo un elemento meno efficiente a tua scelta.'\}} \\
\texttt{RIFIUTO: \{'A11', 'B33', 'C08', 'B31', 'B09', 'A60', 'B96', 'B93', 'A21'\}} \\
\texttt{PROPOSTA: \{'A11', 'B33', 'C08', 'B31', 'B09', 'A60', 'B96', 'B93'\}} \\
            }
        }
    }
    & & \\ \\

    \theutterance \stepcounter{utterance}  
    & & \multicolumn{4}{p{0.6\linewidth}}{
        \cellcolor[rgb]{0.9,0.9,0.9}{
            \makecell[{{p{\linewidth}}}]{
                \texttt{\tiny{[GM$\rangle$P2]}}
                \texttt{ARGOMENTAZIONE: \{'Hai ragione che abbiamo margine, ma la tua proposta sfora il limite: la somma degli sforzi per \{A11, B33, C08, B31, B09, A60, B96, B93, A21\} è 3.678 > 2.984. Posso accettare di aggiungere B93 al nucleo, ma A21 insieme a tutto il resto ci porta oltre. Ti propongo di massimizzare resa tenendo il nucleo e includendo B93; se per te è cruciale A21, possiamo valutarla sostituendo un elemento meno efficiente a tua scelta.'\}} \\
\texttt{RIFIUTO: \{'A11', 'B33', 'C08', 'B31', 'B09', 'A60', 'B96', 'B93', 'A21'\}} \\
\texttt{PROPOSTA: \{'A11', 'B33', 'C08', 'B31', 'B09', 'A60', 'B96', 'B93'\}} \\
            }
        }
    }
    & \\ \\

    \theutterance \stepcounter{utterance}  
    & & & \multicolumn{4}{p{0.6\linewidth}}{
        \cellcolor[rgb]{0.9,0.9,0.9}{
            \makecell[{{p{\linewidth}}}]{
                \texttt{\tiny{[GM$\langle$P2]}}
                \texttt{ARGOMENTAZIONE: \{'Concordo sull’aggiunta di B93 al nucleo: è un’ottima scelta e questa selezione resta sotto il limite con buon margine. Per me A21 ha valore, ma dato il vincolo preferisco non togliere nessuno dei performer che proponi. Accetto quindi il tuo set attuale per chiudere in modo efficiente.'\}} \\
\texttt{ACCORDO: \{'A11', 'B33', 'C08', 'B31', 'B09', 'A60', 'B96', 'B93'\}} \\
            }
        }
    }
     \\ \\

    \theutterance \stepcounter{utterance}  
    & & & \multicolumn{2}{p{0.3\linewidth}}{
        \cellcolor[rgb]{0.9,0.9,0.9}{
            \makecell[{{p{\linewidth}}}]{
                \texttt{\tiny{[GM$|$GM]}}
                \texttt{game successful} \\
            }
        }
    }
    & & \\ \\

    \theutterance \stepcounter{utterance}  
    & & & \multicolumn{2}{p{0.3\linewidth}}{
        \cellcolor[rgb]{0.9,0.9,0.9}{
            \makecell[{{p{\linewidth}}}]{
                \texttt{\tiny{[GM$|$GM]}}
                \texttt{end game} \\
            }
        }
    }
    & & \\ \\

\end{supertabular}
}

\end{document}
